\section*{Hoofdstuk \ref{ch:b3:adaptive-immunity} --- Adaptieve afweer en vaccinatie}

\begin{solution}{exo:b3:adaptive-immunity:1}
Eén \omterm{def:b3:adaptive-immunity:lymphocytes}{lymfocyt}, één specificiteit (Nossal en Lederberg: elke cel scheidde
\omterm{def:b3:adaptive-immunity:antibody}{antilichaam} tegen slechts één \omterm{def:b3:adaptive-immunity:lymphocytes}{antigeen} uit); de receptor wordt gemaakt
voordat het \omterm{def:b3:adaptive-immunity:lymphocytes}{antigeen} wordt ontmoet (Tonegawa: het gen wordt in de \omterm{def:b3:adaptive-immunity:lymphocytes}{B-cel}
herschikt, niet als antwoord op \omterm{def:b3:adaptive-immunity:lymphocytes}{antigeen}); \omterm{def:b3:adaptive-immunity:lymphocytes}{antigeen} selecteert de
passende klonen en breidt ze uit tot effector- en \omterm{def:b3:adaptive-immunity:response}{geheugencellen};
zelfreactieve klonen worden tijdens de ontwikkeling verwijderd.
\end{solution}

\begin{solution}{exo:b3:adaptive-immunity:2}
Twee zware en twee lichte ketens, elk met een variabel domein aan de
punt en constante domeinen erachter; twee identieke bindingsplaatsen,
gevormd door zes hypervariabele lussen; een Fc-stam die andere cellen en
het \omterm{def:b3:innate-immunity:complement}{complement} aflezen. IgM: het eerste \omterm{def:b3:adaptive-immunity:antibody}{antilichaam}, pentameer,
activering van het \omterm{def:b3:innate-immunity:complement}{complement}. IgG: het voornaamste \omterm{def:b3:adaptive-immunity:antibody}{antilichaam} van
bloed en weefsel, \omterm{def:b3:innate-immunity:complement}{opsonisatie}, neutralisatie, passeert de placenta.
IgA: dimeer, uitgescheiden op de slijmvliezen en in de melk. IgE:
gebonden aan mestcellen, tegen parasieten, de bemiddelaar van de
allergie. IgD: op naïeve B-cellen, met weinig bekende functie.
\end{solution}

\begin{solution}{exo:b3:adaptive-immunity:3}
Klasse~I: alle kernhoudende cellen; peptiden uit cytosolische eiwitten,
geknipt door het proteasoom; 8--10 residuen; afgelezen door cytotoxische
CD8-T-cellen. Klasse~II: \omterm{def:b3:innate-immunity:innate}{dendritische cellen}, \omterm{def:b3:innate-immunity:innate}{macrofagen}, B-cellen;
peptiden uit eiwitten die zijn opgenomen en in endosomen verteerd;
13--25 residuen; afgelezen door CD4-helper-T-cellen.
\end{solution}

\begin{solution}{exo:b3:adaptive-immunity:4}
$R_{0}$ is het gemiddelde aantal gevallen dat één geval voortbrengt in
een geheel vatbare bevolking; de drempel is $p_{c} = 1 - 1/R_{0}$. Bof
\qty{80}{\%}; kinkhoest \qty{92}{\%}; griep van 1918 \qty{60}{\%}.
\end{solution}

\begin{solution}{exo:b3:adaptive-immunity:5}
$50\times 20\times 5 = 5000$ zware ketens, $60\times 4 = 240$ lichte:
$1.2\times 10^{6}$ combinaties; met \omterm{def:b3:adaptive-immunity:antibody}{junctionele diversiteit} $1.2\times
10^{10}$ --- tien maal het aantal B-cellen, zodat het dier hoogstens
een tiende tot expressie brengt van wat het zou kunnen maken, en bijna
elke cel een unieke receptor draagt.
\end{solution}

\begin{solution}{exo:b3:adaptive-immunity:6}
$1000\times 10^{-3} = 1$ mutatie per cel per deling; $20$ in twintig
delingen. Onder $10^{5}$ cellen vallen $2\times 10^{6}$
mutatiegebeurtenissen op $3000$ mogelijke enkelvoudige substituties:
elk daarvan wordt honderden malen bemonsterd, en daarbij nog de dubbele
en drievoudige combinaties --- het \omterm{def:b3:adaptive-immunity:germinal-centre}{kiemcentrum} verkent de omgeving van
de receptor uitputtend.
\end{solution}

\begin{solution}{exo:b3:adaptive-immunity:7}
T-cellen worden in de \omterm{def:b3:adaptive-immunity:tolerance}{thymus} geselecteerd om peptiden op de eigen
MHC-allelen van de gastheer te herkennen; de allelen van een ander mens
hebben anders gevormde groeven, die andere peptiden in een andere stand
vasthouden, zodat hetzelfde virale peptide niet wordt gezien. De vreemde
MHC-moleculen van een transplantaat worden zelf door een groot deel van
de T-cellen herkend als ``eigen MHC met een vreemd peptide'', vandaar de
afstoting. Het polymorfisme blijft bestaan omdat een bevolking met vele
allelen een breder monster van de peptiden van elke ziekteverwekker
presenteert en geen ziekteverwekker aan iedereen ontsnapt;
heterozygoten presenteren het dubbele bereik en zijn in het voordeel.
\end{solution}

\begin{solution}{exo:b3:adaptive-immunity:8}
(a) Geen \omterm{def:b3:adaptive-immunity:antibody}{V(D)J-recombinatie}: geen B- of T-cellen, ernstige
gecombineerde immuundeficiëntie. (b) Geen AIRE: cellen van de \omterm{def:b3:adaptive-immunity:tolerance}{thymus}
tonen geen weefselantigenen, zelfreactieve T-cellen ontsnappen,
\omterm{def:b3:adaptive-immunity:tolerance}{auto-immuniteit} tegen vele organen. (c) Geen \omterm{def:b3:adaptive-immunity:tolerance}{FoxP3}: geen regulerende
T-cellen, dodelijke \omterm{def:b3:adaptive-immunity:tolerance}{auto-immuniteit} en allergie op jonge leeftijd.
(d) Geen AID: geen \omterm{def:b3:adaptive-immunity:antibody}{klassenwisseling} en geen hypermutatie --- veel IgM,
bijna geen IgG of IgA, antilichamen met lage affiniteit, terugkerende
besmettingen. (e) Geen CD4-T-cellen: geen hulp voor B-cellen of
\omterm{def:b3:innate-immunity:innate}{macrofagen}, zwakke antilichaamreacties en vatbaarheid voor
opportunistische besmettingen --- de afweerstoornis van aids.
\end{solution}

\begin{solution}{exo:b3:adaptive-immunity:9}
$p_{c} = 1 - 1/6 = 0.83$; dekkingsgraad $0.83/0.9 = 93\,\%$. Bij een
dekkingsgraad van \qty{80}{\%} is het immune deel $0.72$ en
$R_{\text{eff}} = 6\times 0.28
= 1.7$. Onder de vatbaren verspreidt de epidemie zich alsof $R_{0}$
gelijk was aan $1.7$ (de immunen verdunnen alleen de contacten), en de
relatie voor de einduitkomst $s_{\infty} - \ln s_{\infty}/1.7 = 1$
geeft $s_{\infty}
\approx 0.31$: ongeveer \qty{69}{\%} van de vatbaren, \qty{19}{\%} van
de bevolking, wordt besmet.
\end{solution}

\begin{solution}{exo:b3:adaptive-immunity:10}
Uit $\mathrm{d}s/\mathrm{d}t = -\beta si$ en $\mathrm{d}i/\mathrm{d}t
= \beta si - \gamma i$ volgt $\mathrm{d}i/\mathrm{d}s = -1 + \gamma/(\beta s)$,
dus blijft $i + s - (\ln s)/R_{0}$ behouden; met $s = 1$, $i = 0$ aan
het begin en $i = 0$, $s = s_{\infty}$ aan het eind, $s_{\infty} - \ln
s_{\infty}/R_{0} = 1$. Oplossingen: $R_{0} = 1.5$, $s_{\infty} = 0.42$
(\qty{58}{\%} besmet); $2$, $0.20$ (\qty{80}{\%}); $3$, $0.06$
(\qty{94}{\%}); $5$, $0.007$ (\qty{99.3}{\%}). Niet iedereen: naarmate
de vatbaren opraken zakt het effectieve reproductiegetal onder één
terwijl er nog vatbaren over zijn, en dooft de epidemie uit voordat zij
die bereikt.
\end{solution}

\begin{solution}{exo:b3:adaptive-immunity:11}
$1.5^{n} = 1000$: $n = \ln 1000/\ln 1.5 \approx 17$ opeenvolgende
verbeteringen. Met één mutatie per cel per deling en één op honderd die
verbetert, moeten er per ronde minstens honderd cellen worden
doorzocht om er één te vinden; in de praktijk duizenden, want de helft
van de mutaties vernietigt de binding en de selectie is rommelig.
Zeventien ronden van elk twee delingen duren ongeveer negen dagen ---
de waargenomen twee tot drie weken, gegeven de inefficiënties. De cycli
scheiden de mutatie (donkere zone) van de toetsing (lichte zone), zodat
elke variant tegen \omterm{def:b3:adaptive-immunity:lymphocytes}{antigeen} wordt beoordeeld voordat hij zich mag
vermenigvuldigen, precies zoals een genetisch algoritme variatie en
selectie afwisselt.
\end{solution}

\begin{solution}{exo:b3:adaptive-immunity:12}
Vrouwen: verscheidene afweergenen liggen op het X-chromosoom, sommige
(TLR7) ontsnappen aan de inactivering en komen vanaf beide kopieën tot
expressie, en oestrogenen bevorderen de overleving van B-cellen ---
toets: mannen met een extra X (Klinefelter) lopen een vrouwelijk risico
op lupus, en dat blijkt zo te zijn. \omterm{def:b3:adaptive-immunity:mhc}{HLA}: de betrokken allelen
presenteren de betreffende eigen peptiden goed, of presenteren ze niet
in de \omterm{def:b3:adaptive-immunity:tolerance}{thymus}, zodat de klonen niet worden verwijderd --- toets: muizen
die \omterm{def:b3:genetic-engineering:transgenic}{transgeen} zijn voor het menselijke allel krijgen de ziekte wanneer
zij het \omterm{def:b3:adaptive-immunity:lymphocytes}{antigeen} toegediend krijgen. De stijging: minder vroege
blootstelling aan microben en veranderde microbiomen (\omterm{def:b3:bacteriology:antibiotics}{antibiotica},
voeding, keizersnede) geven minder regulerende T-cellen en meer
\omterm{def:b3:innate-immunity:inflammation}{ontsteking} --- toetsen: muizen met aanleg voor diabetes die kiemvrij of
op \omterm{def:b3:bacteriology:antibiotics}{antibiotica} worden grootgebracht krijgen vaker diabetes, en kinderen
van immigranten nemen binnen één generatie de incidentie van het nieuwe
land over.
\end{solution}

\begin{solution}{pb:b3:adaptive-immunity:1}
\textbf{1.} $6000\times 320 = 1.9\times 10^{6}$; $\times 3\times 10^{4}
= 5.8\times 10^{10}$.
\textbf{2.} $10^{11}$ cellen tegenover $5.8\times 10^{10}$ mogelijke:
een lichaam zou elke receptor ongeveer tweemaal kunnen bevatten, maar
de klonen zijn ongelijk en de meeste sequenties worden nooit gemaakt,
zodat elk lichaam een grote maar onvolledige steekproef tot expressie
brengt --- en twee mensen delen slechts een klein deel van hun
receptoren.
\textbf{3.} $10^{11}/10^{5} = 10^{6}$ cellen; in een knoop van
$10^{8}$ ongeveer $1000$.
\textbf{4.} $10^{9}/10^{5} = 10^{4}$ voorlopers --- honderd maal
minder; de pasgeborene heeft bovendien onrijpe follikels en
\omterm{def:b3:innate-immunity:innate}{dendritische cellen}, geen geheugen en beperkte hulp van T-cellen, zodat
de reacties traag en klein zijn, en het IgG van de moeder overbrugt de
kloof.
\textbf{5.} Hij overspant de V--D- en D--J-naden, waar willekeurig
nucleotiden worden verwijderd en toegevoegd: zijn sequentie staat
nergens in het genoom gecodeerd. Omdat hij tussen de twee ketens in het
midden van de bindingsplaats ligt, is hij de lus die het vaakst in
contact met het epitoop staat.
\textbf{6.} Vóór het \omterm{def:b3:adaptive-immunity:lymphocytes}{antigeen} moet de diversiteit breed zijn, zodat
iets bindt wat er ook komt; na het \omterm{def:b3:adaptive-immunity:lymphocytes}{antigeen} is de moeite beter besteed
aan het variëren van de weinige receptoren waarvan al bekend is dat zij
binden. Een grove zoektocht gevolgd door een fijne is veel doelmatiger
dan elk van beide alleen.
\textbf{7.} Zeven dagen zijn $21$ delingen: $100\times 2^{21} = 2.1\times
10^{8}$ cellen; \qty{30}{\%} daarvan, $6.3\times 10^{7}$ \omterm{def:b3:adaptive-immunity:response}{plasmacellen}.
\textbf{8.} $6.3\times 10^{7}\times 2000\times \num{86400} = 1.1\times
10^{16}$ moleculen per dag; $\times 1.5\times 10^{5}\times 1.66\times
10^{-24}$ g $= \qty{2.7}{mg}$; in \qty{3}{L} is dat \qty{0.9}{mg/L} per
dag.
\textbf{9.} \qty{27}{mg} in \qty{3}{L}: \qty{9}{mg/L} $= \qty{0.009}{g/L}$,
een duizendste van het totale IgG --- het \omterm{def:b3:adaptive-immunity:antibody}{antilichaam} tegen één bepaald
\omterm{def:b3:adaptive-immunity:lymphocytes}{antigeen} is een klein deel van het geheel.
\textbf{10.} Honderdvoudig is $\log_{2}100 = 6.6$ halveringstijden,
ongeveer $140$ dagen. \omterm{def:b3:adaptive-immunity:antibody}{Antilichaam} dat jaren aanhoudt komt van
langlevende \omterm{def:b3:adaptive-immunity:response}{plasmacellen} in het merg, die voortdurend uitscheiden, en
van \omterm{def:b3:adaptive-immunity:response}{geheugencellen} die bij hernieuwde blootstelling opnieuw worden
geprikkeld.
\textbf{11.} Negen delingen: $10^{5}\times 512 = 5\times 10^{7}$ cellen
op dag 3 --- duizend maal de $5\times 10^{4}$ van de primaire reactie
op dag 3, en een kwart van wat de primaire reactie pas op dag 7
bereikte.
\textbf{12.} Het herschikte VDJ ligt naast C$_{\mu}$, dus is het eerste
transcript IgM; wisselen naar C$_{\gamma}$ vergt hulp van T-cellen en
AID, en dat kost dagen. \omterm{def:b3:adaptive-immunity:response}{Geheugencellen} zijn al gewisseld, dus maken hun
nakomelingen onmiddellijk IgG.
\textbf{13.} $1000\times 10^{-3} = 1$ per deling; $20$ over twintig.
\textbf{14.} Eén op honderd; $10^{4}\times 0.01 = 100$ verbeterde
cellen per deling.
\textbf{15.} $1.5^{n} = 1000$: $n \approx 17$; $17\times \qty{12}{h}
\approx 8.5$ dagen.
\textbf{16.} De helft van de dochtercellen verliest de binding: zonder
selectie zou na twintig delingen nog maar $0.5^{20}
\approx 10^{-6}$ van de lijn binden. De lichte zone moet de verliezers
bij elke cyclus verwijderen, zodat alleen werkzame en verbeterde
receptoren verdergaan.
\textbf{17.} De eerste dosis laat geheugen-B-cellen achter met
verhoogde affiniteit en gewisselde klasse; de tweede dosis zaait
daarmee nieuwe kiemcentra, en hypermutatie en selectie hervatten vanaf
een hoger beginpunt, zodat de antilichamen die tevoorschijn komen nog
beter zijn.
\textbf{18.} Sommige besmette mensen maken na jaren van rijping
uiteindelijk antilichamen tegen geconserveerde plekken die het \omterm{def:b3:virology:virus}{virus}
niet kan veranderen (breed neutraliserende antilichamen). Het
\omterm{def:b3:adaptive-immunity:germinal-centre}{kiemcentrum} laat zich sturen: een reeks immunogenen die eerst de juiste
kiembaanvoorlopers aanspreekt en daarna de binding aan de geconserveerde
plek beloont, zou de rijping in maanden in plaats van jaren langs die
weg kunnen leiden.
\textbf{19.} $p_{c} = 1 - 1/15 = 93.3\,\%$. Eén dosis: $0.933/0.93 =
100.4\,\%$ --- onbereikbaar; twee doses: $0.933/0.97 = 96.2\,\%$.
\textbf{20.} Immuun $0.90\times 0.97 = 0.873$; $R_{\text{eff}} = 15\times
0.127 = 1.9$: de mazelen circuleren.
\textbf{21.} $60\,000\times (0.10 + 0.90\times 0.03) = 7600$ vatbaren
per geboortecohort; na vijf jaar $38\,000$ --- een reservoir waarin een
ingevoerd geval genoeg vatbaren onder zijn contacten vindt om ketens in
stand te houden.
\textbf{22.} Onder de vatbaren geeft $s_{\infty} - \ln s_{\infty}/1.9 =
1$ de waarde $s_{\infty} \approx 0.23$: ongeveer \qty{77}{\%} van hen,
zo'n $29\,000$ gevallen.
\textbf{23.} Immuun $0.95\times 0.97 = 0.92$; $R_{\text{eff}} = 15\times
0.08 = 1.2$ --- nog steeds boven één; de drempel vergt \qty{96}{\%} met
twee doses, en mazelen is de ziekte die de laatste paar procent
afstraft.
\textbf{24.} Zuigelingen worden alleen beschermd door de immuniteit van
hun omgeving (en enkele maanden door het \omterm{def:b3:adaptive-immunity:antibody}{antilichaam} van de moeder):
kunnen er geen overdrachtsketens ontstaan, dan bereikt geen enkel geval
hen. Bij een dekkingsgraad van \qty{85}{\%} is het immune deel $0.82$ en
$R_{\text{eff}} = 2.6$: er lopen uitbraken, en de zuigelingen, niet
gevaccineerd en het kwetsbaarst, worden besmet.
\textbf{25.} Repertoire ongeveer $6\times 10^{10}$; \qty{2.7}{mg} IgG
per dag van de \omterm{def:b3:adaptive-immunity:response}{plasmacellen} van één afweerreactie; \qty{96}{\%}
dekkingsgraad met twee doses.
\end{solution}
